\section{Discussion}

The central message of this paper is a change of target rather than a new way to summarize persistence. Existing methods answer a question about the conditional law of the diagram given the group label; the estimand developed here asks how the group changes the topology after the covariate distribution is standardized. The two questions are linked in one direction: under randomization $H_0^{\mathrm{cond}}$ implies $H_0^{\mathrm{out}}$, and the converse fails in general, so outside randomization neither implication holds without assumptions that are rarely checked in practice. The target-relative protocol formalizes the consequence: a rejection rate is interpretable only against the null a test actually evaluates, so power comparisons across different nulls are not efficiency statements and cannot be used to rank procedures. The protocol therefore requires that the target-matched class be defined before any procedure is called better than another. The leakage theorems show both failure directions. Under exact causal nullity with covariate shift, $H_0^{\mathrm{cond}}$ fails while $H_0^{\mathrm{out}}$ holds, and every consistent $H_0^{\mathrm{cond}}$ test rejects with probability approaching one; conversely, there are exact laws in which $H_0^{\mathrm{cond}}$ holds while the standardized contrast is nonzero, so no level-$\alpha$ $H_0^{\mathrm{cond}}$ test has power above $\alpha$. These results explain why a conventional benchmark conducted under imbalance cannot be read as evidence about a covariate-adjusted effect.

For practice, the implication is to choose the test from the target and the design, not from a global ranking. If the arm-specific conditional law is the scientific target, or if the design balances covariates by construction, existing tests remain appropriate and adjustment is unnecessary. Adjustment is not a free improvement: the standardized estimand requires causal consistency, exchangeability, and positivity, and the adjusted tests pay a finite-sample price in variance. If covariates drive both membership and topology, the standardized contrast should be estimated with cross-fitted doubly robust scores \cite{chernozhukov_dml_2018,kennedy_drlearner_2023} and tested with one of the two statistics developed here. The outcome-level statistic is the natural default when the scale of the difference is unknown; the distribution-level statistic is preferable when the expected measure itself is the target, and it powered earliest among the distribution-level readouts in the bake-off (a descriptive statement, not a cross-null ranking), but its weak-null calibration requires at least 500 units under strong confounding. The decision table is organized by these conditions: each row names the null tested and the design regime, and no row claims uniform superiority. The protocol also implies a reporting rule: every comparison should print the null it evaluates, because under confounding competitor rates mix signal with size distortion and a bare ranking is not meaningful.

The construction is complementary to existing topological two-sample methodology rather than a competitor on its own turf. Permutation tests on pairwise distances \cite{robinson_turner_2017}, kernel maximum mean discrepancy on persistence diagrams \cite{kwitt_neurips_2015,gretton_2012}, vectorized summaries \cite{moon_lazar_2023,islambekov_pathirana_2024}, Fr\'echet analysis of variance \cite{dubey_muller_2019}, persistence-intensity comparisons \cite{han_kim_kim_2026}, survival-based vectorizations \cite{murris_strand_2026}, and multiscale relevant-difference tests \cite{krebs_rademacher_2024} target either the conditional null or a full distributional equality. Those targets remain legitimate; what changes here is the estimand and the calibration under confounding. The multiple-testing component follows studentized topological multiple testing \cite{vejdemo_mukherjee_2022}, but is built for the grid coordinates of the projected expected measure. The identification requirements also connect the construction to recent work on topological effects under ignorability \cite{kim_lee_tce_2026,souto_diamantis_tcda_2026,saki_faghihi_2026}: standardization is what turns a descriptive contrast into a target with a causal reading, and without the identifying assumptions the estimand drops to the within-stratum contrast rather than the marginal standardized one.

The practical takeaway can be stated compactly. Name the target first. Use an unadjusted method when the conditional law is what the science asks, or when balance makes the distinction immaterial. Use the standardized outcome test when a covariate-adjusted effect is the target. Use the distribution-level test when the expected measure is the target and the sample is large. Treat adjusted tests at small samples with caution, because at $n=50$ the adjusted outcome test is late relative to unadjusted randomization-based competitors while the distribution-level test remains strong. The decision table is a routing device for these choices, not a leaderboard.

\section{Limitations}

Several qualifications bound what this paper establishes. The evidence is synthetic throughout: every rejection rate, power curve, and calibration statement comes from simulated point clouds with known ground truth; there is no real-data analysis, and no application is promised. The distribution-level results have finite-grid scope. The target is the expected persistence measure projected onto one frozen $32 \times 32$ grid, and equality of the unprojected measures is not certified because the total-variation topology on persistence-diagram measures is not separable; the validity, local-limit, and detection-boundary results are fixed-$q$ statements. The framework applies to replicated-cloud designs with one cloud per unit and does not cover the single-cloud two-sample regime, in which each arm contributes one cloud, nor individual-level subgroup, heterogeneity, or conditional-effect questions.

Calibration is the most consequential limitation. Under the strongest confounding cell, the distribution-level weak-null size is 0.096 (Monte Carlo standard error 0.013) at $n=200$, outside the pre-registered $[0.03,0.08]$ band, and 0.066 at $n=500$; we therefore recommend at least 500 units under strong confounding and claim no uniform small-sample size control. Frozen-label permutation is retained only as a diagnostic: it rejects 0.20 to 0.25 under the confounded weak null \cite{chung_romano_2013} and is not finite-sample exact when nuisances are estimated. On multiplicity, the pooled familywise error rates are 0.0392 and 0.0448 over 6000 replications for the two mechanisms, with one inherited conservative miss at 0.028 and one finite-sample size miss at 0.0364 (pooled 0.0312); the miss falsifies finite-sample equality of size, not asymptotic upper control, and no finite-sample exactness is claimed. The studentized comparator is an adaptation of the original procedure, and the multiplier calibration relies on high-dimensional Gaussian approximation \cite{cck_2013}.

The local-power program is reported with its recorded mismatch. The pre-registered 0.06-tolerance screen at the headline cell returned a gap of $+0.102$ (Monte Carlo standard error 0.026) against the naive Gaussian-shift prediction, so the screen was not confirmed; across the 51 scored cells, 9 exceed the 0.06 tolerance and 17 exceed two standard errors, with a maximum gap of 0.275. A finite-$n$ correction that multiplies the out-of-span truncation by an estimated propensity mean ratio moves the headline gap to $-0.053$ on 190 to 300 replications, but its mean-ratio rate is not proved and the completed-record validator still reports the screen as not passed. Both numbers and their criteria are reported, and the bump-direction check was within noise and is not claimed.

Efficiency results are estimator-level at fixed grid dimension, not power optimality. The efficient influence function and the semiparametric bound are proved for the standardized functional and the finite-grid contrast, with attainment at fixed grid dimension, but four of seven pre-registered numerical checks show finite-sample discrepancies and a propensity-related inflation of order $100/n$ that is measured rather than derived, so no finite-sample efficiency claim follows. Uniform double-machine-learning and separation conditions remain assumptions, and the detection-boundary lower-bound clause was not proved. The sharp minimax constant is a conjecture, with only the Le Cam bound and rate slope proved \cite{le_cam_1986,ingster_1993}; no dominance over the persistence-intensity test of Han et al. \cite{han_kim_kim_2026} is claimed, and max-$t$ power monotonicity under positive-semidefinite covariance ordering has 200 of 200 probe support but remains a conjecture. Asymptotic relative efficiency equals one at zero imbalance and the adjusted estimator dominates within the target-matched class under stated conditions, but this is an oracle-nuisance limit; in the completed exact-pairing record the adjusted test sits below the unadjusted test at $n=200$ and $n=500$ in every direction, with mean power gaps of 0.14 to 0.19, so the efficiency record is not numerically confirmed.

The bake-off carries named losses that must be read with the protocol in mind. At $n=50$ under randomization on a mean shift, the adjusted outcome test rejects with probability 0.110 while the unadjusted permutation test, the distance-based test, and the persistence-intensity test reject at 0.800, 1.000, and 1.000. Under strong imbalance in the higher-replication cells, the adjusted outcome test reaches 0.212 (standard error 0.018) at $n=50$ against 0.900 (standard error 0.013) for the distance-based test, and 0.332 (standard error 0.021) at $n=100$ against 0.996 and 0.992 for the distance-based and persistence-intensity tests. The confounded comparisons mix signal with size distortion and are not target-matched rankings; the randomized comparison is the cleaner one and shows a genuine small-sample loss. Coverage is also incomplete: the coarse $n=500$ cells for the loop-driven cloud families and the high-replication mean-shift cell at $n=200$ were not run, on measured cost grounds (about 284 seconds per replication in sequence, and more than two hours for the latter), so the only $n=500$ bake-off evidence is the deterministic separation pilot and no $n=500$ loops-cloud operating point is claimed.

Finally, the leakage program's own gate is inconclusive rather than passed: the false-positive and masking arms are met, but the prototype's size of 0.005 to 0.015 fell outside the pre-registered $[0.03,0.08]$ band, a conservative defect that was superseded by the cross-fitted test. The foundational covariate-shift analysis also required corrections carried through this paper. The loop-count gap at full imbalance is $|0.373|$ (exact quadrature 0.3727, simulation $-0.3738 \pm 0.0013$), not about 1.2; the masking silhouette contrast is $\sup_t \abs{\psi_d(t)} = 0.4934$ (batch standard error $2 \times 10^{-4}$), not 1.45; and a diagnostic formerly labelled as about 0.04 standard errors from zero is about 1.1 standard errors. The masking power statement is now the bound ``at most $\alpha$'' with equality only at a least favorable null, the consistency theorem conditions on labels and excludes covariate-using tests, and the positivity leg is numerical rather than proved. The residual gaps limit the formal completeness of the masking construction but do not affect the headline mechanisms.

\section{Conclusion}

We have argued that covariate adjustment in topological two-sample testing is a question about targets, not about which statistic is sharper. Recasting the comparison as a covariate-standardized contrast under causal consistency, arm-specific exchangeability, and positivity yields two tests built from cross-fitted doubly robust scores: an outcome-level sup-statistic on power-weighted silhouettes and a distribution-level studentized max-$t$ statistic on the projected expected persistence measure. The theory explains when unadjusted tests fail, in both the false-positive and masking directions; gives a local limit, an efficiency bound, and a target-matched efficiency characterization; and certifies fixed-grid validity and shared-multiplier error control under stated conditions. The synthetic evidence shows the adjusted tests broadly holding their level where unadjusted tests do not and often leading the bake-off, while documenting small-sample losses. The honest qualifications are part of the contribution: the distribution-level target is a fixed-grid projection, its calibration is recommended only at $n \ge 500$ under strong confounding, the local-power correction rests on an unproved mean-ratio rate, the efficiency results are estimator-level rather than power-optimal, the sharp minimax constant and max-$t$ monotonicity are conjectural, and the evidence is synthetic. Within those boundaries, the practical message is simple: state the topological target, check the identification assumptions, match the test to the target, and report the null that each comparison evaluates.
